\subsubsection{Ablation on Pre-Calibration}
\label{sec:ablation_precalibration}

\begin{wraptable}[8]{r}{0.40\linewidth}
    \vspace{-1.2em}
    \centering
    \caption{Pre-calibration ablation.}
    \label{tab:ablation_precalibration}
    \scriptsize
    \setlength{\tabcolsep}{8pt}
    \renewcommand{\arraystretch}{1.08}
    \begin{tabular}{lcc}
        \toprule
        \textbf{Metric} & \textbf{Pre-calib.} & \textbf{Dynamic} \\
        \midrule
        Ratio & $1.324\times$ & $1.324\times$ \\
        Escape rate & 0.16\% & 0.16\% \\
        Encode(GB/s) & 430.1 & 80.7 \\
        Decode(GB/s) & 1745.1 & 1779.1 \\
        \bottomrule
    \end{tabular}
    \vspace{-1.5em}
\end{wraptable}

We compare SplitZip's offline pre-calibrated codebook with a dynamic Top-16 variant that rebuilds the exponent codebook for each input.
As shown in Table~\ref{tab:ablation_precalibration}, dynamic calibration achieves the same compression ratio and escape rate as the pre-calibrated design, and decode throughput is almost unchanged.
However, it substantially slows down online compression: encode throughput drops from 430.1~GB/s to 80.7~GB/s, $5.3\times$ slower, because the dynamic variant adds an online histogram and top-$k$ selection pass.
